% !TeX root = ../Book1.tex
\begin{figure}[htbp]
	\centering
	\begin{tikzpicture}[font=\small,scale=0.95]
		\begin{scope}
			\fill[GeometryBlue!10] (0,0)--(30:2.25) arc[start angle=30,end angle=90,radius=2.25]--cycle;
			\foreach \a in {30,90,150}{
				\draw[GeometryGray,dashed] (\a:2.4)--(\a:-2.4);
			}
			\foreach \a in {0,60,120,180,240,300}{
				\draw[->,very thick,GeometryBlue] (0,0)--(\a:1.7);
			}
			\node[right] at (0:1.75) {\(\alpha\)};
			\node[above left] at (120:1.75) {\(\beta\)};
			\node[above right] at (60:1.75) {\(\alpha+\beta\)};
			\node at (60:0.8) {\(C\)};
			\node at (0,-2.85) {\(A_2\)：六根，六个室};
		\end{scope}
		\begin{scope}[xshift=6cm]
			\fill[GeometryBlue!10] (0,0)--(0:2.25) arc[start angle=0,end angle=45,radius=2.25]--cycle;
			\foreach \a in {0,45,90,135}{
				\draw[GeometryGray,dashed] (\a:2.4)--(\a:-2.4);
			}
			\foreach \a in {0,90,180,270}{
				\draw[->,very thick,GeometryBlue] (0,0)--(\a:1.35);
			}
			\foreach \x/\y in {1.35/1.35,-1.35/1.35,-1.35/-1.35,1.35/-1.35}{
				\draw[->,very thick,GeometryBlue] (0,0)--(\x,\y);
			}
			\node[below] at (1.45,0) {\(\varepsilon_1\)};
			\node[above left] at (0,1.4) {\(\varepsilon_2\)};
			\node[right] at (1.35,1.35) {\(\varepsilon_1+\varepsilon_2\)};
			\node at (22.5:0.9) {\(C\)};
			\node at (0,-2.85) {\(B_2\)：八根，八个室};
		\end{scope}
	\end{tikzpicture}
	\caption[秩二根系的根与反射墙]{实根借内积画为蓝色向量，灰色虚线为根的核，浅蓝色区域各标出一个 Weyl 室。左图的根等长；右图的长根长度为短根的 \(\sqrt2\) 倍。图示的是 \(\mathfrak t\) 与对偶空间经内积识别后的超平面排列，不是 \(\mathfrak g\) 中的实根平面。}
	\label{g1:fig:root-arrangements}
\end{figure}
